\pagenumbering{roman}
\setcounter{page}{2}

\chapter*{AGRADECIMIENTOS}
\addcontentsline{toc}{chapter}{AGRADECIMIENTOS}

{
\itshape

A mi familia, por el apoyo sostenido a lo largo de toda la carrera.

A los profesores de la Facultad de Ingeniería y Ciencias Aplicadas de la Universidad de los Andes,
por la formación recibida durante estos años, y de manera especial al profesor guía de este trabajo,
Matías Recabarren, por su orientación y sus observaciones durante el desarrollo del proyecto.

A quienes participaron en las etapas previas de Arbocensus, por construir el proyecto que esta
memoria continúa.
}

\newpage

\tableofcontents
\newpage

% \phantomsection antes de cada \addcontentsline: sin él las tres entradas
% comparten el ancla del capítulo sin numerar anterior y sus enlaces del PDF
% apuntan todos a AGRADECIMIENTOS.
\phantomsection
\addcontentsline{toc}{chapter}{ÍNDICE DE TABLAS}
\listoftables
\newpage

\phantomsection
\addcontentsline{toc}{chapter}{ÍNDICE DE ILUSTRACIONES}
\listoffigures
\newpage

\chapter*{RESUMEN}
\addcontentsline{toc}{chapter}{RESUMEN}

Arbocensus es un proyecto de la Universidad de los Andes orientado a automatizar la gestión de
censos de árboles urbanos. Sus iteraciones anteriores desarrollaron una plataforma web de censado
con participación ciudadana, ludificación y un modelo de estimación de cobertura. El propósito del
proyecto exige un historial de observaciones por árbol, y construirlo obliga a volver sobre un
conjunto conocido de ejemplares ya georreferenciados. Planificar ese re-censo es un problema de
asignación de trabajo que las iteraciones anteriores no abordaron: repartir el arbolado entre
censistas y ordenar sus visitas mediante rutas eficientes, bajo jornadas acotadas.

Esta memoria desarrolla esa etapa mediante una plataforma web que cubre el ciclo operacional
completo del re-censo y un motor de ruteo que lo planifica. El problema se formula como un problema
de ruteo de vehículos de rutas abiertas con visitas opcionales penalizadas y restricción de
duración por ruta ---con el número de rutas, la jornada acotada y el balance de carga resueltos por
el modelo y no declarados por el usuario---, y se resuelve mediante el \textit{solver} de ruteo
vehicular de \textit{OR-Tools}. Los tiempos de desplazamiento peatonal se calculan sobre la red
vial real de \textit{OpenStreetMap} mediante OSRM, lo que elimina costos de APIs externas. El
administrador importa los árboles desde las bases legadas del proyecto, configura los parámetros,
optimiza y publica el plan; el censista accede desde su teléfono a la ruta asignada, con las
paradas en secuencia y navegación peatonal por parada, sobre un despliegue productivo con HTTPS y
entrega continua.

Durante el desarrollo del proyecto se realizaron experimentos sobre doce instancias reales
congeladas de entre $43$ y $1\,607$ árboles, organizados en dieciséis ciclos de evaluación con
criterio de decisión fijado antes de cada corrida y datos crudos versionados. Los ciclos
compararon configuraciones del motor entre sí y contra un heurístico \textit{greedy} de vecino más
cercano sometido al mismo refinamiento local, y doce de ellos cerraron descartando la modificación
que evaluaban, lo que acota el margen de mejora atribuible a la función objetivo.

En la instancia de mayor tamaño, la configuración de producción cubre los 1\,607 árboles con 25
rutas, un índice de balance de carga de 0{,}83 y ninguna ruta que exceda la jornada máxima, en un
presupuesto de cómputo de 120\,s que captura prácticamente toda la mejora alcanzable con
presupuestos mayores, y con un 4{,}7\,\% menos de desplazamiento total ---cerca de 48 minutos---
que el heurístico. En las áreas de censo individuales, la escala de una jornada real, el motor
cubre todos los árboles con conteos bajos de auto-cruces, aunque el contraste frente al heurístico
se cumple en dos de las tres áreas exigidas. Se concluye que la formulación propuesta resuelve el
problema planteado en tiempos compatibles con la planificación de una campaña; su validación en
una jornada de terreno queda pendiente y es la extensión de mayor valor para el proyecto.

\newpage